
\chapter{Unsupported tools}

This Appendix to the ITU-T Software Tool Library (STL) Manual describes the unsupported tools provided in the ITU-T STL.
The tools are named as ``unsupported'' because they do not follow the initial modularity concept of STL.
These tools are provided ``as is'' and without any warranties or implied suitability to use.
However, any feedback on problems with these tools will be welcome and accomodated as possible, as will any improvements made which can be shared and incorporated in the STL.


\section{Source code}
%~~~~~~~~~~~~~~~~~~~~~

\begin{Descr}{40mm}
\item[asc2bin.c:]
        converts decimal or hex ASCII data into short/long/float
                or double binary numbers. Input data must be one
                number per line.

\item[astrip.c:]
        strips off a segment of a file. Can operate on block or
                sample-based parameters and can apply windowing to the
                borders of the extracted segment. Tested in Unix/MSDOS.

\item[bin2asc.c:]
        converts short/long/float or double binary numbers into
                octal, decimal or hex ASCII numbers, printing one per
                line. For Unix/MSDOS.

\item[compfile.c:]
        compare word-wise binary files. For VMS/Unix/MSDOS.

\item[dumpfile.c:]
        dump a binary file. For VMS/Unix/MSDOS.

\item[chr2sh.c:]
           convert {\tt char}-oriented files to \short-oriented (16-bit
           words) files by padding the upper byte of each word of
           the output file with zeros. For Unix/MSDOS.

\item[endian.c:]
  verify whether the current platform is big or little endian
  (i.e. high-byte first or low-byte first). For Unix/MSDOS.

\item[fdelay.c:]
        flexibly introduce delay into a file. Delay
        can be specified in value and length, or can be taken
        from a user-specified file. For Unix/MSDOS.

\item[g728-vt:]
            a directory with software tools for use with the G.728
            floating-point verification package. Not all tools are
            functional; preserved here for future reference.

\item[getcrc32.c:]
         32-bit CRC calculation function and program (depending on
         how it is compiled). Uses the same polynomial as
         ZIP. Checked for portability across a number of
         platforms. Makefile compiles it into an executable called
         crc. For Unix/MSDOS.

\item[measure.c:]
        measure statistics/CRC for a bunch of files.  For
        VMS/Unix/MSDOS.

\item[oper.c:]
            implement arithmetic operation on two files: add,
            subtract, multiply or divide two files applying
            scaling factors (linear or dB), and adding a DC
            level. For Unix/MSDOS.

\item[random.c:]
        randomization tool for selecting items from a list or
            drawing numbers from a range. Reuses the EID linear
            congruential generator from eid.c. For Unix/MSDOS.

\item[sb.c:]
        swap bytes for word-oriented files. For
        VMS/Unix/MSDOS.

\item[sh2chr.c:]
           convert \short-oriented (16-bit words) files to {\tt
           char}-oriented files by ignoring the upper byte of each
           word of the input file. For Unix/MSDOS.

\item[sine.c:]
        generate a sinewave file for a given speco of AC / DC / phase / frequency / sampling frequency values.
        For VMS/Unix/MSDOS.

\item[signal-diff.c:]
        subtract/add files (depending on the compilation, see makefiles). For VMS/Unix/MSDOS.

%\item[\textcolor{blue}{stereoop.c}] \textcolor{blue}{Basic
%        stereo operations on 16-bit audio files. The functionalities
%        include: }

\item[g729e\_convert\_synch.c:]
        convert G.192 word oriented analysis 
        frames without SYNCH\_WORD into a G.192 word oriented 
        file with SYNCH\_WORD and zeroed payload bits. 
         
\end{Descr}


\section{Test files}
%~~~~~~~~~~~~~~~~~~~

\begin{Descr}{40mm}
\item[unsup/test\_data]
Test files for testing some of the unsupported tools: \\
    {\tt\begin{verbatim}
   cf:
     3200    cftest1.dat
     3200    cftest2.dat
     3200    cftest3.dat
      186    delay-15.ref
      186    delay-a.ref
      214    delay-u.ref
      186    delaydft.ref
      200    delayfil.ref

   sb:
      100    bigend.src
      100    litend.src

        \end{verbatim}}
\end{Descr}


%----------------------------------------------------------------------
\section{The {\tt random} tool}
\label{sec:random}
%----------------------------------------------------------------------

%----------------------------------------------------------------------
\subsection{Introduction}
%----------------------------------------------------------------------

The {\tt random} tool provides deterministic pseudo-random selection
capabilities for use in speech and audio codec evaluation scripts. It
was originally developed for the EVS (Enhanced Voice Services) codec
processing scripts as specified in 3GPP S4-121078.

The tool supports two modes of operation:
\begin{itemize}
\item {\bf Subset selection:} randomly select one or more items from a
  provided list, sampling without replacement.
\item {\bf Range mode:} draw random integers from a specified numeric
  range.
\end{itemize}

Deterministic behavior is achieved through a seed-based linear
congruential sequence (LCS) generator, ensuring reproducible results
across platforms when the same seed is used.

%----------------------------------------------------------------------
\subsection{Description of the algorithm}
%----------------------------------------------------------------------

The pseudo-random number generator is inherited from the EID tool
({\tt eid.c}) in the STL. It implements a linear congruential
generator as described by Knuth~\cite{Knuth}:

\[
  \mathrm{seed}_{n+1} = (69069 \times \mathrm{seed}_n + 1) \bmod 2^{32}
\]

The generator returns a uniformly distributed value in $[0, 1)$ by
scaling:
\[
  r = 2^{-32} \times \mathrm{seed}
\]

The {\tt unsigned int} type (32 bits on all modern platforms) is used
to ensure identical behavior on Windows, Linux, and macOS.

\subsubsection{Subset mode}

In subset mode, the tool selects {\em n} items from a list of {\em N}
items without replacement. For each selection:
\begin{enumerate}
\item A random index $k = \lfloor N \cdot r \rfloor$ is computed.
\item The item at position $k$ is output.
\item The item is removed from the list (remaining items shift down).
\item $N$ is decremented for the next draw.
\end{enumerate}

\subsubsection{Range mode}

In range mode, the tool draws {\em n} random integers from the
inclusive range $[\mathrm{start}, \mathrm{stop}]$. Each value is computed
as:
\[
  v = \mathrm{start} + \lfloor (\mathrm{stop} - \mathrm{start} + 1) \cdot r \rfloor
\]

Note that in range mode, values are drawn independently (with
replacement).

\subsubsection{Dummy pre-runs}

An optional number of dummy pre-runs can be specified. These advance
the RNG state before selection begins, effectively providing a
secondary diversification mechanism beyond the seed value.

%----------------------------------------------------------------------
\subsection{Implementation}
%----------------------------------------------------------------------

The tool is implemented as a single source file {\tt random.c} in the
{\tt src/unsup/} directory.

\subsubsection{Usage}

{\tt
\begin{verbatim}
random [OPTIONS] [ITEM_LIST]
\end{verbatim}
}

\subsubsection{Options}

\begin{Descr}{40mm}
\item[\tt -s SEED]
        Seed for the random number generator. Any value between 0 and
        $2^{32}-1$. Default: 3141592653.

\item[\tt -d PRERUNS]
        Number of dummy pre-runs to advance the RNG state before
        selection begins. Default: 0.

\item[\tt -r START STOP]
        Range mode: draw random integers from the inclusive range
        $[\mathrm{START}, \mathrm{STOP}]$. When this option is specified,
        no item list is expected.

\item[\tt -n NUM\_ITEMS]
        Number of items to select (subset mode) or numbers to draw
        (range mode). Default: 1.

\item[\tt -v]
        Verbose output: shows the internal state of the item list
        after each selection (subset mode only).

\item[\tt -h]
        Print usage information and exit.
\end{Descr}

\subsubsection{Examples}

\paragraph{Select one item from a list:}
{\tt
\begin{verbatim}
$ random -s 42 alpha bravo charlie delta echo
alpha
\end{verbatim}
}

\paragraph{Select 3 items from a list with 3 dummy pre-runs:}
{\tt
\begin{verbatim}
$ random -s 42 -d 3 -n 3 alpha bravo charlie delta echo
bravo delta charlie
\end{verbatim}
}

\paragraph{Draw 5 random integers from a range:}
{\tt
\begin{verbatim}
$ random -s 12345 -r 10 20 -n 5
12 19 12 14 14
\end{verbatim}
}

\subsubsection{Limitations}

\begin{itemize}
\item Maximum of 1000 items in subset mode.
\item The seed is parsed with {\tt atoi()}, limiting practical input to
  signed integer range on some platforms. For full 32-bit range,
  numeric overflow wraps as expected for unsigned arithmetic.
\end{itemize}

%----------------------------------------------------------------------
\subsection{Tests and portability}
%----------------------------------------------------------------------

Three CTest test cases validate the tool:

\begin{enumerate}
\item {\bf random1:} Range mode with seed 12345, range $[10, 20]$, 5
  draws. Expected output: ``12 19 12 14 14''.
\item {\bf random2:} Subset mode with seed 42, selecting 1 item from 5.
  Expected output starts with ``alpha''.
\item {\bf random3:} Subset mode with seed 42, 3 dummy pre-runs,
  selecting 3 items from 5. Expected output: ``bravo delta charlie''.
\end{enumerate}

Tests use {\tt PASS\_REGULAR\_EXPRESSION} to verify stdout matches
expected values, ensuring cross-platform reproducibility of the RNG
sequence.

